% RESEARCH_DRAFT_SHA256: b1e03eafdcbb10ec0b95c33e48fb8d1db73268e70104e2d619ee8f08e863a31d
% Layout and source-reference wrappers expanded for the web renderer.
\documentclass[11pt]{article}
\usepackage[margin=25mm]{geometry}
\usepackage{fontspec,unicode-math}
\setmainfont{Latin Modern Roman}
\setmathfont{Latin Modern Math}
\usepackage{amsmath,amssymb,mathtools,xurl,hyperref}
\hypersetup{colorlinks=true,urlcolor=blue,linkcolor=blue}
\setcounter{secnumdepth}{0}
\setlength{\parindent}{0pt}
\setlength{\parskip}{5pt}
\setlength{\emergencystretch}{3em}
\newcommand{\sref}[2]{\hyperlink{source:#1:#2}{[S#2]}}
\newcommand{\cataloguescope}[1]{\paragraph{Source record.}#1}
\newcommand{\floor}[1]{\left\lfloor#1\right\rfloor}
\begin{document}
\section*{Research-attempt provenance}
This standalone research attempt retains the complete original definition body below,
including its historical source and preparation notes. Those notes describe the
earlier catalogue preparation; the new attempt and its verification are identified
separately after that body. No original ID, catalogue status, or source statement is changed.

\section{Scaling Paradigms for Lean Autoformalization}\label{statement}

\subsection{Definitions and mathematical statement}

\paragraph{Shared notation.}Standard number systems are $\mathbb N,\mathbb Z,\mathbb Q,\mathbb R,\mathbb C$. Specialized notation and index ranges follow the cited source and the definitions below.


A proof-search tree branches over permitted proof steps. Autoformalization constructs a formal theorem statement and proof from another representation; debugging revises candidates after formal checking fails.
\paragraph{Record scope and qualifications.}
The catalogue title is retained as metadata. Named certificates, contracts, bounds or frameworks in that title are not asserted as results of the original source extract.
The task distribution, resource budget, semantic-correctness standard, and success metric are not supplied.

Source-record limitation: An empirical comparison of theorem-proving paradigms is requested without a formal computational target.
\paragraph{Statement recorded in the supplied source.}
Determine which combination of fine-grained tree search, whole-proof generation, iterative debugging, hierarchical decomposition, and reinforcement learning scales most effectively to large research-level theorem-proving and autoformalization tasks in Lean. \href{https://doi.org/10.4171/owr/2025/43}{[S2]}
\paragraph{Source status.}
The supplied catalogue labels this source record open; that metadata was checked against the supplied snapshot on 2026-10-11. The label is preserved for provenance and is not an independent verification of present open status or mathematical correctness. The catalogue literature assessment is dated 2026-08-22.
\subsection{Short English statement}
Compare ways to scale formal proof search and autoformalization.
\subsection{Sources}
\begin{itemize}
\item[S1] Ulam AI and the UnsolvedMath Contributors, UnsolvedMath: A Curated Collection of Open Mathematics Problems, catalogue entry. Catalogue curation and metadata: CC BY 4.0. The source extract is retained; definitions, English summaries and qualifications are editorial additions.
\url{https://huggingface.co/datasets/ulamai/UnsolvedMath}.
\item[S2] MATRIX-MFO Tandem Workshop: Machine Learning and AI for Mathematics (2026), Oberwolfach Reports, DOI 10.4171/owr/2025/43, pp. 2323--2350
\url{https://doi.org/10.4171/owr/2025/43}.
\item[S3] Improving Autoformalization using Type Checking
\url{https://arxiv.org/abs/2406.07222}.
\item Source-specific Oberwolfach Reports licences, including CC BY-SA 4.0 where applicable, are retained; see THIRD\_PARTY\_NOTICES.md.
\end{itemize}
\paragraph{Preparation and review.}
The original source and UnsolvedMath contributors retain attribution for the source extract. Codex (GPT-6; exact serving revision not exposed) assisted with definitions, English summaries, scope notes and formatting. No new human mathematical authorship or independently verified solution is asserted. Catalogue identity checked against the complete supplied registry and current repository. Distinct upstream records, including repeated questions and broad or damaged extracts, are retained for complete coverage of the source CS category. A repeated record is not a new independent conjecture. Complete mathematical-equivalence and current-literature review has not been performed. Preparation and final review follow contribution prompts C and D, with source gaps disclosed. Structural and rendering checks do not establish a new theorem or open status.
\label{end:statement}

\section{Literature and context check}
The primary report, printed p. 2347, compares fine-grained search, whole-proof generation, iterative debugging, and hierarchical approaches and asks which scales best for Lean reinforcement learning. No task distribution, trained policy, verifier cost, or compute budget is fixed by the question. The following exact search models show why a scaling conclusion needs those choices rather than merely the names of the paradigms.
\section{Research attempt: Two search-oracle models with opposite scaling conclusions}

Fix an alphabet of size $b\geq2$ and a target proof of length $d\geq1$. In the first model, every correct prefix has exactly one valid next symbol, and a local oracle immediately recognizes whether a proposed next symbol is the valid one. A uniformly random whole string is correct with probability $b^{-d}$. Independent complete trials therefore have expected count $b^d$: summing the geometric tail gives $\sum_{k\geq0}(1-b^{-d})^k=b^d$. Generating every trial in full costs expected $d b^d$ symbols.

A prefix-preserving method tries the next symbols without replacement in a uniformly random order, retaining a correct prefix. The valid symbol's position is uniform on $\{1,\ldots,b\}$, so the expected number of tests per level is $(b+1)/2$ and the total is $d(b+1)/2$. This proves an exponential advantage of local feedback under these very informative oracle assumptions, with cost measured in symbol tests.

Now change only the feedback model. Every prefix of length less than $d$ is locally valid and receives the same feedback, while exactly one of the $b^d$ complete strings is successful. Before a success, the only information is that tested complete strings are wrong. If the hidden successful string is uniform, any deterministic adaptive method testing distinct strings encounters it at a uniform position in its query order: after each failure all unqueried strings remain symmetric. Thus its expected complete-string tests are $(b^d+1)/2$. Randomized methods satisfy the same bound by conditioning on their randomness; repeats can only waste queries. Prefix caching may save repeated token work but cannot produce informative pruning in this model.

These are not claims about the distribution of Lean proof states. They expose the exact missing variable: how much a local success or failure predicts eventual completion. In actual proof search, a locally valid step can leave an unprovable or prohibitively hard goal; it need not be the unique step on a successful proof. Conversely a strong whole-proof generator is not uniformly sampling strings, so its success probability need not be $b^{-d}$. Hierarchical decomposition also changes the verifier and proposal distributions and cannot be compared without charging those costs.

A fair empirical or theoretical scaling statement must therefore specify at least a family of goals, proposal laws, reuse of verified prefixes or lemmas, and a common resource measure. The two proved models refute an attempt to deduce a universally superior paradigm solely from the fact that one uses finer steps.
\section{Verification and adversarial audit}
Expected costs are derived from the exact query distributions; they do not compare symbol tests to full-proof tests without stating the difference. The uninformative-oracle lower bound allows adaptive order and randomized algorithms. The first model explicitly assumes immediate perfect local feedback, which is stronger than ordinary acceptance of a proof-state transition.
\paragraph{Audit corrections.} Separated local validity from eventual solvability and made the proposal distribution and counted resource explicit in each scaling comparison.
\paragraph{External computation.} None; symbolic arguments and boundary checks only.
\section{Scope and remaining gaps}
No Lean scaling experiment, reinforcement-learning algorithm, or universal paradigm ranking is obtained. The original agenda remains unresolved, with the feedback and cost-model gap stated precisely.
\paragraph{Integration review.} Codex (GPT-6; exact serving revision not exposed) reviewed the written arguments and expanded layout and source-reference wrappers. Repository metadata and the production web renderer were checked using Python 3.13 and JavaScript with MathJax 3.2.2. This model-assisted review is not independent human proof verification; no PDF compilation or formal proof checking is asserted.

\section{Completion Estimate}
\noindent\textbf{COMPLETION: 0\%}
The calculations below the preserved source statement investigate a concrete route,
but establish no substantiated new progress toward the entire original target.
Known special cases and failed sufficient conditions are not credited as a solution.
\section{References}
The original source references and catalogue attribution are preserved in the
Sources section above. The following URLs were actually consulted for this attempt;
retention of an additional original reference does not claim fresh consultation of it.
\begin{itemize}
\item \url{https://ems.press/content/serial-article-files/52253} (consulted 2026-10-11). 
\end{itemize}
\section{Model and contributor disclosure}
Codex GPT-6 completed written-proof review of this attempt at its declared scope; this is model-assisted review, not independent human proof checking.
OpenAI gpt-6-astra, requested reasoning configuration ultra, generated this research
attempt on 2026-10-11. The exact serving revision is not exposed. Codex GPT-6
performed orchestration and repository-format review; its exact serving revision
is also not exposed. The mathematical route, boundary checks and final adversarial
self-audit followed contribution prompts A then B in
\texttt{docs/contribution\_prompts.md}, after reading \texttt{CONTRIBUTING.md}.
Python 3.14.2 assembled the standalone documents; any mathematical computation is
specified in Verification. No theorem prover was used. No human mathematical
authorship, substantial human proof correction, or independent human proof checking
is asserted. Model-assisted review, source retrieval, and structural checks do not
establish correctness or novelty. Full private reasoning transcripts are unavailable
as reproducibility artifacts; the complete arguments, inputs and outcomes actually
used are stated here or in the companion batch audit.
\section{Final status}
Unresolved as a new research contribution. This is a concrete route investigation,
not a claim of new complete resolution or a current-status certification. A known
answer reproduced for context does not become a new contribution. The catalogue's
source-status label is preserved; this attempt status does not assert that a
source-labelled solved question is open.
\end{document}
